\label{sec:introduction}

Every redistricting cycle produces a set of electoral consequences whose political
salience far exceeds their frequency: incumbent pairings. When two sitting members of
Congress find themselves in the same new district, at least one career is disrupted.
The affected legislators face three options: compete in a primary against a fellow
incumbent (expensive, divisive, and uncertain), seek an open seat in a different
district (requiring campaign infrastructure transplantation), or retire. None of these
options is attractive. Incumbent legislators therefore place enormous pressure on state
mapmakers to draw districts that minimise pairing, and mapmakers with access to incumbent
home addresses have consistently accommodated this pressure~\citep{altman2011public}.

The question this paper addresses is simple but consequential: how does this deliberate
pairing-minimisation compare to what a politically blind algorithm would produce?
The \textsc{bisect} redistricting system~\citep{bisect2024} draws congressional districts
by recursively bisecting a census tract adjacency graph to optimise population balance,
with optional edge weights derived from geographic characteristics. The algorithm receives
no information about incumbent home addresses. It cannot minimise pairing because it does
not know where incumbents live.

This creates a natural experiment. We observe the pairing rate under two plans---the
\textsc{bisect} algorithmic map and the enacted 2022 congressional map---for the same
set of incumbents in the same geographic state. The difference between these rates is
attributable to exactly one factor: the enacted mapmakers' access to and use of incumbent
address data. If the algorithmic map produces pairing rates consistent with the random
baseline, and the enacted map produces pairing rates significantly below that baseline,
we have strong evidence that enacted pairing protection is a deliberate redistricting
choice requiring incumbency data, not a consequence of geographic constraints.

\subsection{Contribution}

This paper makes four contributions:

\begin{enumerate}
  \item \textbf{Analytical baseline}: We derive the exact expectation and variance for
    the incumbent pairing count under uniform random redistricting for NC, WI, and TX,
    and provide a geographic adjustment that accounts for the spatial clustering of
    incumbent residences.
  \item \textbf{Algorithmic benchmark}: We report the pairing rate under \textsc{bisect}
    algorithmic maps for all three states and test whether these rates are consistent
    with the random baseline.
  \item \textbf{Enacted map comparison}: We test whether enacted 2022 pairing rates
    differ significantly from the random baseline, establishing the magnitude of deliberate
    incumbency protection.
  \item \textbf{Legal interpretation}: We interpret these findings in the context of
    redistricting litigation, identifying the conditions under which pairing-rate
    comparisons can serve as evidence of strategic mapmaking.
\end{enumerate}

\subsection{States Studied}

We study three states representing distinct redistricting contexts. North Carolina
($k=14$) is a historically contested redistricting state with a long record of federal
litigation over both racial and partisan gerrymandering~\citep{shaw1993, harper2022}.
Wisconsin ($k=8$) provides a smaller-$k$ context where pairing probabilities are
lower but the enacted map's zero-pairing outcome is more conspicuous against the random
baseline. Texas ($k=38$) provides the largest-$k$ case where even a random process
would produce substantial pairing, making the enacted map's near-zero outcome more
analytically remarkable.

All results use 2020 Census tract data and FEC 2022 candidate filing addresses
geocoded to census tract centroids. Single-seed \textsc{bisect} results are marked
with $^\dagger$.
